\documentclass[10pt,a4paper]{amsart}
\usepackage[margin=19mm]{geometry}
\usepackage{amsmath,amssymb,mathtools}
\usepackage[scr=rsfs,cal=euler]{mathalfa}
\usepackage{xfrac}
\usepackage[unicode,hidelinks,bookmarks=false]{hyperref}
\newcommand{\ie}{i.e.\ }
\newcommand{\cf}{cf.\ }
\newcommand{\resp}{respectively\ }
\newcommand{\Subsection}{\subsection}
\providecommand{\nobreakdash}{\nobreak\hbox{-}}
\setlength{\emergencystretch}{2em}
\multlinegap0pt
\usepackage{amssymb}
\usepackage{enumerate}
\usepackage{bm}
\usepackage{cancel}
\newtheorem{lemma}{Lemma}
\newtheorem{theorem}{Theorem}
\newtheorem{proposition}{Proposition}
\newcommand{\id}{\mathrm{id}}
\newcommand{\qo}{{\boldsymbol{\preccurlyeq}}}
\newcommand{\qop}{{\boldsymbol{\succcurlyeq}}}
\newcommand*{\vcbox}[1]{\begingroup
\setbox0=\hbox{#1}\parbox{\wd0}{\box0}\endgroup}
\title{Monoidal quantaloids}
\author{Gejza Jen\v{c}a, Bert Lindenhovius}
\date{Source: arXiv:2504.18266v3 (CC BY 4.0)}
\begin{document}
\maketitle

\noindent\emph{Source and licence.} Monoidal quantaloids. Version arXiv:2504.18266v3 (CC BY 4.0), distributed by arXiv under the Creative Commons Attribution 4.0 International licence, http://creativecommons.org/licenses/by/4.0/, as recorded in the arXiv licence metadata for this exact version and shown on the abstract page https://arxiv.org/abs/2504.18266v3. Full licence notice: \texttt{licenses/arxiv-2504.18266-CC-BY-4.0.txt}.\par\smallskip

\noindent\textit{Editorial scope.} Counted unit: the article's theorem theorem (source0.tex lines 2522--2524). The article's complete original proof of it is delivered with it (source0.tex lines 2526--2580, one \texttt{proof} environment of 55 lines). No mathematics is written, repaired or re-derived: the statement and the proof are the article's own lines, in the article's own notation, and no retained line is substituted or re-flowed.  Retained uncounted as the source-local context the proof needs: lines 541--549; lines 2319--2326; lines 2498--2507.  One declared adaptation: exactly 20 ancillary strings inside the retained spans are omitted (image-inclusion commands and, where the source repeats a label, the repeated label definitions) and no other line differs from the source; the figure environments, with their placement, captions and labels, are kept, and the package records the omitted strings in fidelity receipts bound to the affected bodies.  No retained line contains a citation, so the wrapper carries no bibliography and no bibliography entry.  The retained lines contain the article's own diagram code, which the wrapper typesets with the standard tikz packages; no image file is delivered and the package declares no asset dependency.  Editorial formatting only, in addition: an AMS \texttt{amsart} preamble that restates the article's own declarations and macros used by the retained lines, namely the environments \texttt{lemma}, \texttt{theorem}, \texttt{proposition} on separate counters, together with the standard packages those lines need.  The retained environments print in this document as Lemma 1, Theorem 1, Proposition 1 in order of appearance, whereas the article prints the same items with the numbering of its own full document.  The complete native source of the article as distributed in the arXiv e-print for this exact version is bundled unchanged as \texttt{sources/177081/source0.tex}.
\begin{lemma}\label{lem:dagger_and_dual}
    Let $f:X\to Y$ be a morphism in a dagger compact category $(\mathbf C,\otimes, I)$. 
    Then 
    \begin{align*}
        (f^*)^\dag & =(f^\dag)^*;\\
        \epsilon_Y\circ(f\otimes \id_{Y^*}) & = \epsilon_X\circ(\id_X\otimes f^*);\\
        (\id_{X^*}\otimes f)\circ \eta_X & =(f^*\otimes \id_Y)\circ \eta_Y.
          \end{align*}
\end{lemma}

\begin{theorem}\label{thm:monoidal structure PreOrd}
Let $(\mathbf R,\otimes,I)$ be a dagger symmetric monoidal quantaloid. The category $\mathrm{PreOrd}(\mathbf R)$ becomes a symmetric monoidal category as follows:
\begin{itemize}
    \item We define the monoidal product by $(X,\qo_X)\otimes(Y,\qo_Y):=(X\otimes Y,\qo_X\otimes\qo_Y)$ on objects, and on monotone maps by the monoidal product of their underlying morphisms in $\mathbf R$;
    \item The monoidal unit is $(I, \id_I)$;
    \item the associator, unitors and symmetry between preordered objects are the respective associator, unitors and symmetry between the underlying objects of $\mathbf R$.
\end{itemize}   Moreover, the inclusion functor $J:\mathrm{Maps}(\mathbf R)\to\mathrm{PreOrd}(\mathbf{R})$, $X\mapsto (X,\id_X)$ is strict monoidal, and left adjoint to the forgetful functor $U:\mathrm{PreOrd}(\mathbf R)\to\mathrm{Maps}(\mathbf R)$, $(X,\qo)\mapsto X$.
\end{theorem}

\begin{proposition}\label{prop:MonRel is symmetric monoidal}
   Let $(\mathbf R,\otimes,I)$ be a dagger symmetric monoidal quantaloid. Then
   $\mathrm{MonRel}(\mathbf R)$ is a symmetric monoidal quantaloid if we equip it with a monoidal product $\otimes$ as follows:
   \begin{itemize}
       \item The monoidal product $\otimes$ coincides with the monoidal product on $\mathrm{PreOrd}(\mathbf R)$ as defined in Theorem \ref{thm:monoidal structure PreOrd}, i.e., $(X,\qo_X)\otimes (Y,\qo_Y)=(X\otimes Y,\qo_X\otimes\qo_Y)$;
       \item the monoidal product $r\otimes s$ of monotone relations $r$ and $s$ is given by the monoidal product of $r$ and $s$ in $\mathbf R$;
       \item the monoidal unit is given by $(I,\id_I)$;
       \item if $\alpha$, $\lambda$, $\rho$ and $\sigma$ denote the respective associator, left unitor, right unitor and symmetry of $(\mathrm{PreOrd}(\mathbf R),\otimes,(I,\id_I))$, then  the associator, left unitor, right unitor and symmetry of $\mathrm{MonRel}(\mathbf R)$ are given by $\alpha_\diamond$, $\lambda_\diamond$, $\rho_\diamond$, and $\sigma_\diamond$, respectively.
   \end{itemize}
\end{proposition}

\begin{theorem}
    Let $(\mathbf R,\otimes,I)$ be a dagger compact quantaloid with respective unit and counit morphisms $\eta_X:I\to X^*\otimes X$ and $\epsilon_X:X\otimes X^*\to I$ for each object $X$. Then $(\mathrm{MonRel}(\mathbf R),\otimes,(I,\id_I))$ is a compact quantaloid with respective unit and counit morphism $\eta_{(X,\qo)}:(I,\id_I)\to (X,\qo)^*\otimes(X,\qo)$ and $\epsilon_{(X,\qo)}:(X,\qo)\otimes(X,\qo)^*\to(I,\id_I)$ given by $\eta_{(X,\qo)}:=(\qop^*\otimes\qop)\circ\eta_X$ and $\epsilon_{(X,\qo)}:=\epsilon_X\circ(\qop\otimes\qop^*)$.
\end{theorem}

\begin{proof}
By Proposition \ref{prop:MonRel is symmetric monoidal}, $\mathrm{MonRel}(\mathbf R)$ is a symmetric monoidal quantaloid. 

Let $(X,\qo)$ be a preordered object of $\mathbf R$. Since $\qo\circ\qo\leq \qo$, and $\id_X\leq\qo$, we have $\qo=\id_X\circ\qo\leq\qo\circ\qo\leq\qo$, whence $\qo\circ\qo=\qo$, so preorders are idempotent. We will use this in the remainder of proof without mentioning it. Note that by Lemma \ref{lem:dagger_and_dual}, we have $(\qo^*)^\dag=(\qo^\dag)^*=\qop^*$. Note furthermore that if $f:(X,\qo)\to(Y,\qo_Y)$ is an order isomorphism between preordered objects, functoriality of $(-)_\diamond$ yields that $f_\diamond$ is also invertible in $\mathrm{MonRel}(\mathbf R)$, and $(f_\diamond)^{-1}=(f^{-1})_\diamond=\qop\circ f$. We simply write $f_\diamond^{-1}$ instead of  $(f_\diamond)^{-1}$ or $(f^{-1})_\diamond$.
We have show to that the  unit and counit satisfy the zigzag identities of a compact closed category, for which we
used string diagrams. By definition, we have the following identities.
\begin{center}
\vcbox{}
\vcbox{$\overset{(D1)}{=}$}
\vcbox{}
\qquad\qquad
\vcbox{}
\vcbox{$\overset{(D2)}{=}$}
\vcbox{}
\qquad\qquad
\vcbox{}
\vcbox{$\overset{(D3)}{=}$}
\vcbox{}
\end{center}
Furthermore, we have the following identities, where the first is the first zigzag identity for $\mathbf R$ as a compact category, the second by Lemma \ref{lem:dagger_and_dual}, and the third by combining reflexivity and transitivity in the definition of a preorder, respectively:
\begin{center}
\vcbox{}
\vcbox{$\overset{(I1)}{=}$}
\vcbox{}
\qquad\qquad
\vcbox{}
\vcbox{$\overset{(I2)}{=}$}
\vcbox{}
\qquad\qquad
\vcbox{}
\vcbox{$\overset{(I3)}{=}$}
\vcbox{}
\end{center}
Then, suppressing associators and unitors, as usual with string diagrams, we find:
\begin{center}
\vcbox{}
\vcbox{$\overset{(D1)-(D3)}{=}$}
\vcbox{}
\vcbox{$\overset{(I3)}{=}$}
\vcbox{}
\vcbox{$\overset{(I2)}{=}$}\\
\end{center}
\begin{center}
\vcbox{}
\vcbox{$\overset{(I3)}{=}$}
\vcbox{}
\vcbox{$\overset{(I1)}{=}$}
\vcbox{}
\vcbox{$\overset{(I3)}{=}$}
\vcbox{}
\vcbox{$\overset{(D3)}{=}$}
\vcbox{ }
\end{center}
In a similar way, we show that $\mathrm{MonRel}(\mathbf R)$ satisfy the second zigzag identity of compact-closed categories, which proves the statement.
\end{proof}

\end{document}
